%%=============================================================================
%% Validatie & Testing
%%=============================================================================

\chapter{\IfLanguageName{dutch}{Validatie en testing}{Validation and testing}}%
\label{ch:validatie}

Dit hoofdstuk beschrijft het testplan en de testresultaten waarmee de correctheid, veiligheid en schaalbaarheid van de PoC worden aangetoond.

%%----------------------------------------------------------------------
\section{Testplan}%
\label{sec:validatie-testplan}

\subsection{Doelstellingen}

% TODO: Beschrijf wat het testplan beoogt te valideren:
% - Werkt het dynamisch aanmaken en beëindigen van containers correct?
% - Is de isolatie tussen instanties gegarandeerd?
% - Zijn de beveiligingsmaatregelen effectief?
% - Houdt het systeem stand onder zware belasting?

\subsection{Testomgeving}

% TODO: Beschrijf de testomgeving:
% - Hardware/VM-specificaties
% - Betrokken software en versies (CTFd, Docker, plugin, OS)
% - Netwerktopologie van de testopstelling

%%----------------------------------------------------------------------
\section{Functionele tests}%
\label{sec:validatie-functioneel}

\subsection{Aanmaken en beëindigen van containers}

% TODO: Beschrijf de testscenario's:
% - Deelnemer start een challenge: wordt de container correct aangemaakt?
% - Deelnemer stopt de challenge: wordt de container correct verwijderd?
% - Timeout: wordt de container automatisch gestopt na de ingestelde tijd?
% - Meerdere deelnemers starten tegelijk dezelfde challenge.

\subsection{Isolatie tussen gebruikers}

% TODO: Beschrijf hoe wordt geverifieerd dat containers van verschillende
% deelnemers elkaar niet kunnen bereiken:
% - Netwerktest: kan container A container B pingen?
% - Kan een deelnemer de omgeving van een andere deelnemer manipuleren?

%%----------------------------------------------------------------------
\section{Beveiligingstests}%
\label{sec:validatie-beveiliging}

\subsection{Container escape pogingen}

% TODO: Beschrijf welke aanvalstechnieken werden getest:
% - Misbruik van gemounte volumes
% - Privilege escalation via capabilities
% - Docker socket exposure
% Welke beschermingen verhinderden succesvolle escapes?

\subsection{Resource exhaustion simulaties}

% TODO: Beschrijf hoe resource exhaustion werd gesimuleerd:
% - Fork bomb binnen een container
% - Geheugenoverbelasting
% Hoe reageerde het systeem? Werden de resource-limieten afgedwongen?

%%----------------------------------------------------------------------
\section{Load testing}%
\label{sec:validatie-loadtest}

\subsection{Testopzet}

% TODO: Beschrijf de load test:
% - Gebruikte tool: k6 \autocite{k6} of Locust \autocite{Locust}
% - Testscenario: X gelijktijdige gebruikers starten een challenge
% - Duur en ramp-up profiel

\subsection{Meetpunten}

% TODO: Welke metrics werden gemeten?
% - Latency bij het aanmaken van een container
% - CPU- en RAM-gebruik op de host(s)
% - Foutpercentage (mislukte container-starts)
% - Doorvoer (containers per seconde)

\subsection{Resultaten en conclusies}

% TODO: Presenteer de testresultaten in tabellen of grafieken.
% Wat zijn de grenzen van de PoC? Bij welk aantal gelijktijdige gebruikers
% treedt er degradatie op?
